\section{Účel systému}
\label{sec:kelvin-purpose}

Hlavním účelem systému Kelvin je podpora výuky programování, především v jazycích C, C++, Python nebo Rust.
Systém poskytuje jednotné webové rozhraní, prostřednictvím kterého mohou studenti odevzdávat svá řešení úloh.
Vyučujícím následně umožňuje tato řešení kontrolovat, poskytovat k nim zpětnou vazbu a udělovat bodové hodnocení.

\subsection{Typický scénář použití}
\label{subsec:kelvin-workflow}

Cyklus jedné úlohy v systému Kelvin probíhá následovně:

Vyučující nejprve vytvoří úlohu: zadá její textový popis, definuje testovací vstupy a očekávané výstupy a nastaví termín odevzdání.
Úloha je pak zpřístupněna přihlášeným studentům okmažitě nebo v předem stanoveném čase.

Student si zadání přečte, vypracuje řešení ve svém vývojovém prostředí a výsledný zdrojový kód odevzdá prostřednictvím webového rozhraní systému.
Bezprostředně po odevzdání systém spustí automatizované testy: zdrojový kód je zkompilován, spuštěn proti připraveným testovacím vstupům a výsledky jsou porovnány s očekávanými výstupy.
Student vidí výsledky testů okamžitě -- které testy prošly a které selhaly -- a může na základě zpětné vazby řešení opravit a odevzdat znovu.

Po uzavření termínu odevzdání přistoupí vyučující k manuálnímu hodnocení.
Prohlíží zdrojový kód odevzdaného řešení přímo v prostředí systému, přidává komentáře ke konkrétním řádkům nebo blokům kódu a na základě celkového posouzení udělí bodové hodnocení.
Student hodnocení a komentáře obdrží a může na ně reagovat, čímž vzniká asynchronní diskuse nad zdrojovým kódem.

Systém Kelvin je rovněž využíván pro testy, kvízy a jiné formy hodnocení, které nemusí nutně zahrnovat zdrojový kód -- například pro testování teoretických znalostí nebo algoritmického myšlení prostřednictvím krátkých textových odpovědí.

\subsection{Řádkové komentáře jako klíčová funkce}
\label{subsec:kelvin-line-comments}

Jednou z nejvýznamnějších funkcí systému z pohledu praxe je možnost přidávat komentáře přímo ke konkrétním řádkům zdrojového kódu.
Vyučující může označit libovolný řádek nebo rozsah řádků a k němu připsat připomínku.
Komentáře jsou studentovi zobrazeny přímo v kontextu kódu, takže okamžitě vidí, ke které konkrétní části řešení se zpětná vazba vztahuje.

Tato vlastnost je z pohledu integrace LLM zcela zásadní, protože LLM generuje komentáře přirozeně ve formátu "k řádku X patří připomínka Y".
Výstup modelu tak lze přímo mapovat na existující datový model systému Kelvin bez nutnosti zavádět nové typy entit.

\endinput